\SourceLeaf{21}
\begin{center}\Large SUMMATION\end{center}

\begin{center}\large OF CERTAIN SPECIAL SERIES.\end{center}

\BookRule{22mm}

\Article{1}

Among the more notable truths to which the theory of circle division has opened the way, the summation proposed in Disquiss. Arithmet., Article 356, holds a distinguished place. This is due both to its particular elegance and remarkable fruitfulness, which another investigation will later give us occasion to discuss more fully, and to the unusual difficulties that beset its rigorous and complete proof.\quad These difficulties should have been all the less expected because they concern a certain restriction of the theorem more than the theorem itself: once this restriction is disregarded, a proof is immediately available and follows very easily from the theory developed in that work.\quad The theorem was given there in the following form.\quad Suppose that $n$ is a prime number. Represent all the quadratic residues of $n$ between $1$ and $n-1$ inclusive by $a$, and all the nonresidues between the same limits by $b$; finally, let $\omega$ be the arc $\frac{360^{0}}{n}$, and $k$ any fixed integer not divisible by $n$. Then we have

I.\quad for a value of $n$ of the form $4m+1$,

\[
\begin{aligned}
\Sigma\cos ak\omega&=-\tfrac12\pm\tfrac12\surd n\\
\Sigma\cos bk\omega&=-\tfrac12\mp\tfrac12\surd n,\ \text{and hence}\\
\Sigma\cos ak\omega-\Sigma\cos bk\omega&=\pm\surd n\\
\Sigma\sin ak\omega&=0\\
\Sigma\sin bk\omega&=0
\end{aligned}
\]

\SourceLeaf{22}
II.\quad for a value of $n$ of the form $4m+3$,

\[
\begin{aligned}
\Sigma\cos ak\omega&=-\tfrac12\\
\Sigma\cos bk\omega&=-\tfrac12\\
\Sigma\sin ak\omega&=\pm\tfrac12\surd n\\
\Sigma\sin bk\omega&=\mp\tfrac12\surd n\\
\Sigma\sin ak\omega-\Sigma\sin bk\omega&=\pm\surd n
\end{aligned}
\]

These summations were proved with full rigor in the place cited, and the only difficulty that remains here is to determine the \emph{sign} to be placed before the radical.\quad It can indeed be shown without difficulty that this sign depends on the number $k$ to the following extent: the \emph{same} sign must hold for all values of $k$ that are quadratic residues of $n$, and the opposite sign for all values of $k$ that are quadratic nonresidues of $n$.\quad Thus the whole question turns on the value $k=1$: as soon as the sign for this value is known, the signs for all the other values of $k$ are immediately available as well.\quad Yet this very question, which at first seems to be among the easier ones, presents us with unforeseen difficulties, and the method that had guided us this far without impediment utterly refuses any further help.

\Article{2}

Before proceeding further, it will be useful to work out some examples of our summation numerically; we should, however, first make a few general observations.

I.\quad In the case where $n$ is a prime number of the form $4m+1$, let all the quadratic residues of $n$ between $1$ and $\frac12(n-1)$ inclusive be represented by $a'$, and all the nonresidues between the same limits by $b'$. All the $n-a'$ are known to be among the $a$, and all the $n-b'$ among the $b$. Since all the $a'$, $b'$, $n-a'$, $n-b'$ together evidently make up the entire collection of numbers $1$, $2$, $3\dotsfour$ $n-1$, all the $a'$ together with all the $n-a'$ comprise all the $a$; likewise all the $b'$ together with all the $n-b'$ comprise all the $b$.\quad Hence

\SourceLeaf{23}
\[
\begin{aligned}
\Sigma\cos ak\omega&=\Sigma\cos a'k\omega+\Sigma\cos(n-a')k\omega\\
\Sigma\cos bk\omega&=\Sigma\cos b'k\omega+\Sigma\cos(n-b')k\omega\\
\Sigma\sin ak\omega&=\Sigma\sin a'k\omega+\Sigma\sin(n-a')k\omega\\
\Sigma\sin bk\omega&=\Sigma\sin b'k\omega+\Sigma\sin(n-b')k\omega
\end{aligned}
\]

Now, since $\cos(n-a')k\omega=\cos a'k\omega$, $\cos(n-b')k\omega=\cos b'k\omega$, $\sin(n-a')k\omega=-\sin a'k\omega$, $\sin(n-b')k\omega=-\sin b'k\omega$, it follows at once that

\[
\begin{aligned}
\Sigma\sin ak\omega&=\Sigma\sin a'k\omega-\Sigma\sin a'k\omega=0\\
\Sigma\sin bk\omega&=\Sigma\sin b'k\omega-\Sigma\sin b'k\omega=0
\end{aligned}
\]

The summation of the cosines takes the form

\[
\begin{aligned}
\Sigma\cos ak\omega&=2\Sigma\cos a'k\omega\\
\Sigma\cos bk\omega&=2\Sigma\cos b'k\omega
\end{aligned}
\]

so that we must have

\[
\begin{aligned}
1+4\Sigma\cos a'k\omega&=\pm\surd n\\
1+4\Sigma\cos b'k\omega&=\mp\surd n\\
2\Sigma\cos a'k\omega-2\Sigma\cos b'k\omega&=\pm\surd n
\end{aligned}
\]

\Needspace{9\baselineskip}II.\quad In the case where $n$ is of the form $4m+3$, the complement of any residue $a$ to $n$ is a nonresidue, and the complement of any $b$ is a residue. Thus all the $n-a$ coincide with all the $b$, and all the $n-b$ with all the $a$. This gives

\[
\begin{aligned}
\Sigma\cos ak\omega=\Sigma\cos(n-b)k\omega=\Sigma\cos bk\omega
\end{aligned}
\]

Since all the $a$ and $b$ together make up all the numbers $1$, $2$, $3\dotsfour n-1$, and hence $\Sigma\cos ak\omega+\Sigma\cos bk\omega=\cos k\omega+\cos2k\omega+\cos3k\omega+\text{etc.}$ $+\cos(n-1)k\omega=-1$, the summations

\[
\begin{aligned}
\Sigma\cos ak\omega&=-\tfrac12\\
\Sigma\cos bk\omega&=-\tfrac12
\end{aligned}
\]

follow immediately.\quad Likewise,

\[
\begin{aligned}
\Sigma\sin ak\omega=\Sigma\sin(n-b)k\omega=-\Sigma\sin bk\omega
\end{aligned}
\]

\SourceLeaf{24}
which shows how the two summations

\[
\begin{aligned}
2\Sigma\sin ak\omega&=\pm\surd n\\
2\Sigma\sin bk\omega&=\mp\surd n
\end{aligned}
\]

depend on one another.

\Article{3}

Here are the numerical calculations for several examples:

I.\quad For $n=5$ there is one value of $a'$, namely $a'=1$, and one value of $b'$, namely $b'=2$. We have

\[
\begin{aligned}
\cos\omega=+0{,}3090169944\qquad\cos2\omega=-0{,}8090169944
\end{aligned}
\]

and hence $1+4\cos\omega=+\surd5$,\quad $1+4\cos2\omega=-\surd5$.

II.\quad For $n=13$ there are three values of $a'$, namely $1$, $3$, $4$, and as many values of $b'$, namely $2$, $5$, $6$. We calculate

{\small\setlength{\arraycolsep}{2.5pt}\renewcommand{\arraystretch}{1.15}
\[
\begin{array}{rclcrcl}
\cos\omega&=&+0{,}8854560257&\qquad&\cos2\omega&=&+0{,}5680647467\\
\cos3\omega&=&+0{,}1205366803&\qquad&\cos5\omega&=&-0{,}7485107482\\
\cos4\omega&=&-0{,}3546048870&\qquad&\cos6\omega&=&-0{,}9709418174\\\cline{1-3}\cline{5-7}
\text{Sum}&=&+0{,}6513878190&\qquad&\text{Sum}&=&-1{,}1513878189
\end{array}
\]
}

Hence $1+4\Sigma\cos a'\omega=+\surd13$,\quad $1+4\Sigma\cos b'\omega=-\surd13$.

\Needspace{13\baselineskip}III.\quad For $n=17$ we have four values of $a'$, namely $1$, $2$, $4$, $8$, and as many values of $b'$, namely $3$, $5$, $6$, $7$.\quad The cosines are therefore calculated as follows:

{\small\setlength{\arraycolsep}{2.5pt}\renewcommand{\arraystretch}{1.15}
\[
\begin{array}{rclcrcl}
\cos\omega&=&+0{,}9324722294&\qquad&\cos3\omega&=&+0{,}4457383558\\
\cos2\omega&=&+0{,}7390089172&\qquad&\cos5\omega&=&-0{,}2736629901\\
\cos4\omega&=&+0{,}0922683595&\qquad&\cos6\omega&=&-0{,}6026346364\\
\cos8\omega&=&-0{,}9829730997&\qquad&\cos7\omega&=&-0{,}8502171357\\\cline{1-3}\cline{5-7}
\text{Sum}&=&+0{,}7807764064&\qquad&\text{Sum}&=&-1{,}2807764065
\end{array}
\]
}

Hence $1+4\Sigma\cos a'\omega=+\surd17$,\quad $1+4\Sigma\cos b'\omega=-\surd17$.

\SourceLeaf{25}
IV.\quad For $n=3$ there is just one value of $a$, namely $a=1$, to which corresponds

\[
\begin{aligned}
\sin\omega=+0{,}8660254038
\end{aligned}
\]

Hence $2\sin\omega=+\surd3$.

V.\quad For $n=7$ there are three values of $a$, namely $1$, $2$, $4$. These give the sines

{\small\setlength{\arraycolsep}{2.5pt}\renewcommand{\arraystretch}{1.15}
\[
\begin{array}{rcll}
\sin\omega&=&+0{,}7818314825&\\
\sin2\omega&=&+0{,}9749279122&\\
\sin4\omega&=&-0{,}4338837391&\\\cline{1-3}
\text{Sum}&=&+1{,}3228756556,&\text{and hence }2\Sigma\sin a\omega=+\surd7.
\end{array}
\]
}

VI.\quad For $n=11$ the values of $a$ are $1$, $3$, $4$, $5$, $9$, to which correspond the sines

{\small\setlength{\arraycolsep}{2.5pt}\renewcommand{\arraystretch}{1.15}
\[
\begin{array}{rcll}
\sin\omega&=&+0{,}5406408175&\\
\sin3\omega&=&+0{,}9898214419&\\
\sin4\omega&=&+0{,}7557495744&\\
\sin5\omega&=&+0{,}2817325568&\\
\sin9\omega&=&-0{,}9096319954&\\\cline{1-3}
\text{Sum}&=&+1{,}6583123952,&\text{and hence }2\Sigma\sin a\omega=+\surd11.
\end{array}
\]
}

\Needspace{19\baselineskip}VII.\quad For $n=19$ the values of $a$ are $1$, $4$, $5$, $6$, $7$, $9$, $11$, $16$, $17$, to which correspond the sines

{\small\setlength{\arraycolsep}{2.5pt}\renewcommand{\arraystretch}{1.15}
\[
\begin{array}{rcll}
\sin\omega&=&+0{,}3246994692&\\
\sin4\omega&=&+0{,}9694002659&\\
\sin5\omega&=&+0{,}9965844930&\\
\sin6\omega&=&+0{,}9157733267&\\
\sin7\omega&=&+0{,}7357239107&\\
\sin9\omega&=&+0{,}1645945903&\\
\sin11\omega&=&-0{,}4759473930&\\
\sin16\omega&=&-0{,}8371664783&\\
\sin17\omega&=&-0{,}6142127127&\\\cline{1-3}
\text{Sum}&=&+2{,}1794494718,&\text{and hence }2\Sigma\sin a\omega=+\surd19.
\end{array}
\]
}

\SourceLeaf{26}
\Article{4}

In all these examples the radical takes the positive sign. This is also easily confirmed for the larger values $n=23$, $n=29$, etc., and already makes it highly probable that the same holds in general.\quad But a proof of this phenomenon cannot be drawn from the principles set out in the place cited, and it must with full justification be regarded as a matter for deeper investigation.\quad The purpose of this paper is therefore to present a rigorous proof of this most elegant theorem, attempted in vain in various ways over several years and finally achieved through unusual and quite subtle considerations. At the same time we shall bring the theorem to a far greater degree of generality, preserving, or rather increasing, its elegance.\quad Finally, as a crowning addition, we shall show the remarkable and very close connection between this summation and another arithmetical theorem of the greatest importance.\quad We hope that these investigations will themselves be welcome to geometers, and that the methods by which we have accomplished all this, and which may also be useful on other occasions, will be considered worthy of their attention.

\Article{5}

Our proof comes from considering a special kind of series whose terms depend on expressions such as

\[
\begin{aligned}
\frac{(1-x^m)\,(1-x^{m-1})\,(1-x^{m-2})\dotsfour(1-x^{m-\mu+1})}{(1-x)\,(1-xx)\,(1-x^3)\dotsfive(1-x^\mu)}
\end{aligned}
\]

For brevity we shall denote such a fraction by $(m,\mu)$, and first make a few general observations about functions of this kind.

\Needspace{19\baselineskip}I.\quad Whenever $m$ is a positive integer less than $\mu$, the function $(m,\mu)$ evidently vanishes, since the numerator contains the factor $1-x^0$.\quad For $m=\mu$, the factors in the numerator are identical, in reverse order, to those in the denominator, so that $(\mu,\mu)=1$. Finally, for the case where $m$ is a positive integer greater than $\mu$, we have the formulas

\[
\begin{aligned}
(\mu+1,\mu)&=\frac{1-x^{\mu+1}}{1-x}=(\mu+1,1)\\
(\mu+2,\mu)&=\frac{(1-x^{\mu+2})\,(1-x^{\mu+1})}{(1-x)\qquad(1-xx)}=(\mu+2,2)\\
(\mu+3,\mu)&=\frac{(1-x^{\mu+3})(1-x^{\mu+2})(1-x^{\mu+1})}{(1-x)\qquad(1-xx)\qquad(1-x^3)}=(\mu+3,3)\ \text{etc.}
\end{aligned}
\]

\SourceLeaf{27}
or, in general,

\[
\begin{aligned}
(m,\mu)=(m,m-\mu)
\end{aligned}
\]

II.\quad Further, it is easily confirmed that in general

\[
\begin{aligned}
(m,\mu+1)=(m-1,\mu+1)+x^{m-\mu-1}(m-1,\mu)
\end{aligned}
\]

and since we likewise have

\[
\begin{aligned}
(m-1,\mu+1)&=(m-2,\mu+1)+x^{m-\mu-2}(m-2,\mu)\\
(m-2,\mu+1)&=(m-3,\mu+1)+x^{m-\mu-3}(m-3,\mu)\\
(m-3,\mu+1)&=(m-4,\mu+1)+x^{m-\mu-4}(m-4,\mu),\ \text{etc.}
\end{aligned}
\]

and can continue this series as far as

\[
\begin{aligned}
(\mu+2,\mu+1)&=(\mu+1,\mu+1)+x(\mu+1,\mu)\\
&=(\mu,\mu)+x(\mu+1,\mu)
\end{aligned}
\]

provided that $m$ is a positive integer greater than $\mu+1$, we obtain

\[
\begin{aligned}
(m,\mu+1)&=(\mu,\mu)+x(\mu+1,\mu)+xx(\mu+2,\mu)\\
&\qquad+x^3(\mu+3,\mu)+\text{etc.}\\
&\qquad+x^{m-\mu-1}(m-1,\mu)
\end{aligned}
\]

\Needspace{16\baselineskip}This shows that if, for some fixed value of $\mu$, every function $(m,\mu)$ is a polynomial whenever $m$ is a positive integer, then every function $(m,\mu+1)$ must also be a polynomial.\quad Since that supposition holds for $\mu=1$, it therefore also holds for $\mu=2$, and hence for $\mu=3$, etc. That is, in general, for any positive integral value of $m$, $(m,\mu)$ is a polynomial, or equivalently the product

\[
\begin{aligned}
(1-x^m)(1-x^{m-1})(1-x^{m-2})\dotsfour(1-x^{m-\mu+1})
\end{aligned}
\]

is divisible by

\[
\begin{aligned}
(1-x)(1-x^2)(1-x^3)\dotsfour(1-x^\mu)
\end{aligned}
\]

\clearpage\Needspace{14\baselineskip}\Article{6}

We shall now consider two series, both of which can lead to our goal.\quad The first series is

\SourceLeaf{28}
\[
\begin{aligned}
&1-\frac{1-x^m}{1-x}+\frac{(1-x^m)\,(1-x^{m-1})}{(1-x)\,(1-xx)}\\
&\qquad-\frac{(1-x^m)(1-x^{m-1})(1-x^{m-2})}{(1-x)\quad(1-xx)\quad(1-x^3)}+\text{etc.}
\end{aligned}
\]

or

\[
\begin{aligned}
1-(m,1)+(m,2)-(m,3)+(m,4)-\text{etc.}
\end{aligned}
\]

which for brevity we shall denote by $f(x,m)$.\quad First, it is immediately evident that whenever $m$ is a positive integer this series \emph{terminates} after its term in position $m+1$ (with $=\pm1$), and that in this case its sum must therefore be a finite polynomial in $x$.\quad Further, Article 5, II, shows that in general, for any value of $m$,

\[
\begin{aligned}
1&=1\\
-(m,1)&=-(m-1,1)-x^{m-1}\\
+(m,2)&=+(m-1,2)+x^{m-2}(m-1,1)\\
-(m,3)&=-(m-1,3)-x^{m-3}(m-1,2),\ \text{etc.}
\end{aligned}
\]

and hence

\[
\begin{aligned}
f(x,m)&=1-x^{m-1}-(1-x^{m-2})(m-1,1)\\
&\qquad+(1-x^{m-3})(m-1,2)\\
&\qquad-(1-x^{m-4})(m-1,3)+\text{etc.}
\end{aligned}
\]

But evidently

\[
\begin{aligned}
(1-x^{m-2})(m-1,1)&=(1-x^{m-1})(m-2,1)\\
(1-x^{m-3})(m-1,2)&=(1-x^{m-1})(m-2,2)\\
(1-x^{m-4})(m-1,3)&=(1-x^{m-1})(m-2,3),\ \text{etc.}
\end{aligned}
\]

from which we deduce the equation

\[
\begin{aligned}
f(x,m)=(1-x^{m-1})f(x,m-2)\ .\ .\ .\ .\ .\ .\ .\ .\ [1]
\end{aligned}
\]

\clearpage\Needspace{11\baselineskip}\Article{7}

Since for $m=0$ we have $f(x,m)=1$, the formula just found gives

\[
\begin{aligned}
f(x,2)&=1-x\\
f(x,4)&=(1-x)(1-x^3)\\
f(x,6)&=(1-x)(1-x^3)(1-x^5)\\
f(x,8)&=(1-x)(1-x^3)(1-x^5)(1-x^7),\ \text{etc.}
\end{aligned}
\]

or, in general, for any even value of $m$,

\[
\begin{aligned}
f(x,m)=(1-x)(1-x^3)(1-x^5)\dotsfour(1-x^{m-1})\ .\ .\ .\ [2].
\end{aligned}
\]

\SourceLeaf{29}
Conversely, since for $m=1$ we have $f(x,m)=0$, we also have

\[
\begin{aligned}
f(x,3)&=0\\
f(x,5)&=0\\
f(x,7)&=0\ \text{etc.}
\end{aligned}
\]

\Needspace{5\baselineskip}or, in general, for any odd value of $m$,

\[
\begin{aligned}
f(x,m)=0
\end{aligned}
\]

The latter summation could, moreover, already have been derived from the fact that in the series

\[
\begin{aligned}
1-(m,1)+(m,2)-(m,3)+\text{etc.}+(m,m-1)-(m,m)
\end{aligned}
\]

the last term cancels the first, the next to last cancels the second, etc.

\Article{8}

For our purpose the case where $m$ is a positive odd integer is sufficient. But the unusual nature of the matter makes it worthwhile to add a few remarks about the cases where $m$ is fractional or negative.\quad Our series then evidently no longer terminates, but continues to infinity; it is also easily seen that it becomes divergent whenever $x$ is given a value less than $1$, so that its summation must be restricted to values of $x$ greater than $1$.

\Needspace{13\baselineskip}By formula [1] of Article 6 we have

\[
\begin{aligned}
f(x,-2)&=\frac{1}{1-\frac1x}\\[8pt]
f(x,-4)&=\frac{1}{1-\frac1x}\pdot\frac{1}{1-\frac1{x^3}}\\[8pt]
f(x,-6)&=\frac{1}{1-\frac1x}\pdot\frac{1}{1-\frac1{x^3}}\pdot\frac{1}{1-\frac1{x^5}},\ \text{etc.}
\end{aligned}
\]

so that the value of the function $f(x,m)$ can also be given in finite terms for a negative even integral value of $m$.\quad For the remaining values of $m$, however, we shall convert the function $f(x,m)$ into an \emph{infinite product} as follows.

As $m$ tends to a negative \emph{infinite} value, the function $f(x,m)$ becomes

\SourceLeaf{30}
\[
\begin{aligned}
&1+\frac1{x-1}+\frac1{x-1}\pdot\frac1{xx-1}\\
&\qquad+\frac1{x-1}\pdot\frac1{xx-1}\pdot\frac1{x^3-1}+\text{etc.}
\end{aligned}
\]

This series is therefore equal to the infinite product

\[
\begin{aligned}
\frac1{1-\frac1x}\pdot\frac1{1-\frac1{x^3}}\pdot\frac1{1-\frac1{x^5}}\pdot\frac1{1-\frac1{x^7}}\;\text{etc. to infinity.}
\end{aligned}
\]

Further, since in general

\[
\begin{aligned}
f(x,m)&=f(x,m-2\lambda)\pdot(1-x^{m-1})(1-x^{m-3})\\
&\qquad(1-x^{m-5})\dotstwo(1-x^{m-2\lambda+1})
\end{aligned}
\]

we have

\[
\begin{aligned}
f(x,m)&=f(x,-\infty)\pdot(1-x^{m-1})(1-x^{m-3})(1-x^{m-5})\\
&\qquad\text{etc. to infinity.}\\
&=\frac{1-x^{m-1}}{1-x^{-1}}\pdot\frac{1-x^{m-3}}{1-x^{-3}}\pdot\frac{1-x^{m-5}}{1-x^{-5}}\pdot\frac{1-x^{m-7}}{1-x^{-7}}\\
&\qquad\text{etc. to infinity.}
\end{aligned}
\]

whose factors evidently converge ever more closely to unity as they continue.

The case $m=-1$ deserves particular attention. Here

\[
\begin{aligned}
f(x,-1)=1+x^{-1}+x^{-3}+x^{-6}+x^{-10}+\text{etc.}
\end{aligned}
\]

\Needspace{9\baselineskip}This series is therefore equal to the infinite product

\[
\begin{aligned}
\frac{1-x^{-2}}{1-x^{-1}}\pdot\frac{1-x^{-4}}{1-x^{-3}}\pdot\frac{1-x^{-6}}{1-x^{-5}}\;\text{etc.}
\end{aligned}
\]

or, writing $x$ for $x^{-1}$, we obtain

\[
\begin{aligned}
1+x+x^3+x^6+\text{etc.}&=\frac{1-xx}{1-x}\pdot\frac{1-x^4}{1-x^3}\pdot\frac{1-x^6}{1-x^5}\\
&\qquad\pdot\frac{1-x^8}{1-x^7}\;\text{etc.}
\end{aligned}
\]

This equality between two rather intricate expressions, to which we shall return on another occasion, is certainly very remarkable.
